\scanpage{17}
% Pencil foliation: 15, upper right; 16, lower left.
c'est un autre foncteur fibre, mais \uline{gradué},
donc défini par un \uline{torseur} $P\deleted{\unsure{G}}$ sur $G$, et
\inserted{$F$ se \unsure{trouve|déduit} de $F'$ par torseur $P'$ sur $G'$.}
Supposons que sur \deleted{\unsure{quelque part}} une extension
$K'$ de $K$, on puisse trouver un isom
$F\sim F'$ compatible avec les filtrations, et
induisant l'identité sur les gradués associés.
[\uuline{Ex} \inserted{$S=\operatorname{Spec}k$,} $k$ de car. nulle,
$F=T_{\mathrm{DR}}$ avec filtration
canonique ; ça marche encore si $S$ est n'importe
quel schéma \inserted{régulier} de car. nulle \dots].
Alors
il y a un sous-torseur \inserted{canonique} $Q$ de $P\deleted{\unsure{G}}$, de
groupe $H$ le sous-groupe de $G$ des autom
compatibles avec filtration et induisant l'identité
sur le gradué, et de \unsure{$=$|$\overline{\mathrm m}$} \deleted{\unsure{pour}} un $Q'$ dans
$P'$, \inserted{groupe $H'$.}
On notera que $H,H'$ sont unipotents. Ainsi,
$F$ se déduit canoniquement du foncteur
gradué $F'$ [\unsure{opération|application} définie par
\diagram{\begin{tikzcd}
\mathbb{G}_{\mathrm m,K}\arrow[r,"\alpha"]&G'
\end{tikzcd}}]
à l'aide d'un torseur $Q'$ sur le groupe
\inserted{unipotent} $H(\alpha)$ de $G'$ défini par $\alpha$ comme on verra.

\margin{\figure{\begin{tikzpicture}
\node[draw,ellipse,align=left] {expliciter\\la construc-\\tion de\\$H(\alpha)$};
\end{tikzpicture}}}

13) \uline{Relations avec le groupe de
\deleted{Galois}\inserted{Mumford-Tate}.}

\inserted{Les structures de Hodge polarisables forment une $\otimes$-catégorie
semi-simple. Tte sous-catégorie de type fini est ``définie'' par un
groupe alg. $G$ sur $\mathbb{Q}$.}

Soit $(V,V_0,(V^{p,q}))$ une structure de Hodge. $^{b}$

\margin{$b$ avec
$i_1,i_2:\mathbb{G}_{\mathrm m,\mathbb{C}}\longrightarrow G_{\mathbb{C}}$
satisfaisant $i_1(\lambda)i_2(\lambda)=i(\lambda)$ central,
$i_2(\lambda)=\overline{i_1(\unsure{\lambda|\overline\lambda})}$.
On a réciproquement}
